\begin{foldable} % Model compression --> Knowledge distillation
    One inevitable challenge encompassing the advancement of artificial intelligence (AI) that often gets overlooked is the sheer size of AI models.
    This obstacle limits their deployability in constrained environments, such as mobile or edge devices.
    Fortunately, research on \textit{ model compression} addresses this problem through various techniques such as model pruning, quantization, or low-rank factorization~\cite{gou2021knowledge}.
    Among these techniques, \textit{knowledge distillation} (KD)~\cite{hinton2015distilling} is well known for its efficiency and intuitive approach.
\end{foldable}

\begin{foldable} % Knowledge distillation
    The intuition of KD bears much resemblance to that of human learning: transferring the knowledge from a complex expert network (\textit{teacher}) to a simple novice network (\textit{student}).
    In conventional machine learning, the student learns to predict the ground-truth labels from a dataset, much like a human learner does self-studying given a collection of exercises and solutions.
    The advantageous ingredient introduced by KD is the additional supervision by a skillful teacher with more informative hints, improving the student's generalization.
    Thus, the student may reach comparable performance to the teacher despite having a lesser capacity.
    This foundational work~\cite{hinton2015distilling} has inspired a populated line of KD research~\cite{romero2015fitnets, chen2017learning, park2019relational, meng2019conditional, nguyen2021knowledge}.
\end{foldable}

\begin{foldable} % Data-free KD
    However, traditional KD methods often require the teacher's original training data.
    Ideally, the knowledge domain covered by the original data should provide a perfect match to the teacher's expertise, facilitating a seamless knowledge transfer process.
    Nonetheless, the source dataset is often protected by stringent regulations as they are tied to security, privacy, and proprietary concerns.
    In fact, personal and health data, or data as intellectual properties all qualify for protection under~\cite{eu2016gdpr, us1996hipaa, eu1996databasedirective, us2016dsta} imposed by the EU and US governments.
    These regulations effectively inhibit the use of traditional KD methods and pose the more challenging \textit{data-free} KD problem.
\end{foldable}

\begin{foldable} % Black-box data-free KD
    Another fundamental requirement for the success of KD is the access to the teacher model, varying from \textit{white-box} (\ie, access to the intermediate features, parameters, activations, or gradients, \etc)~\cite{romero2015fitnets, yim2017gift, chen2019data} to \textit{black-box} (\ie, only the final logits or predictive probabilities)~\cite{wang2020neural, nguyen2022black, vo2024improving}.
    Unfortunately, AI models are often shipped as a closed API, returning only the final post-processed outputs.
    For instance, large language models from \cite{openai2022, anthropic2024} only return textual replies (instead of per-token scores), or certain evaluation servers, \eg, \cite{russakovsky2015imagenet}, only return the top-1 prediction.
    Retrieving deeper insights via these interfaces is usually costly or even non-negotiable.
    In conjunction, the \textit{black-box data-free KD} (BBDFKD) setting invalidate most KD methods and call for a fresh approach.
\end{foldable}

\begin{foldable} % Weakness of existing methods: ZSDB3
    To the best of our knowledge, a handful of methods have attempted to tackle the BBDFKD problem using synthetic images for KD.
    Within these methods, IDEAL~\cite{zhang2022ideal} and DFHL-RS~\cite{yuan2024data} leverage a generator network to generate synthetic images.
    Differently, ZSDB3~\cite{wang2021zero} samples synthetic images from random uniform noise and then updates them as trainable parameters.
    Critically, all three did not leverage \textit{diversity} for distillation.
    IDEAL and DFHL-RS attempt to generate class-balanced and realistic images, but are at risk of mode collapse caused by generator overfitting.
    Whereas in ZSDB3, the synthetic images have a suboptimal initialization strategy, where some classes could be unforgivingly unobtainable.
\end{foldable}

\begin{foldable} % Our method
    \textbf{Our method.} We tackle the BBDFKD problem, \ie, distilling the student from a black-box teacher which only returns top-1 prediction and without access to any real data.
    Focusing on data diversity, we realize our solution with three components: \textit{Synthesis} --- to craft a pipeline of image priors with diverse distinct patterns, \textit{Contrastion} --- to improve their diversity through contrastive learning, and \textit{Distillation} --- to train a high-performance student from these images.
\end{foldable}

\begin{foldable} % Contribution
    \textbf{Contribution.} In summary, we would like to deliver our solution \textit{\textbf{D}iverse \textbf{I}mage \textbf{P}riors \textbf{K}nowledge \textbf{D}istillation} (\textbf{\methodname}) with the following contributions:
    \begin{enumerate}
        \item We propose \textit{image priors} --- synthetic images crafted with distinct and diverse natural semantics.
        \item We further investigate contrastive learning as a fitting strategy to improve image diversity for KD.
        \item We outperform previous BBDFKD methods on multiple benchmarks and conduct comprehensive ablation study.
    \end{enumerate}
\end{foldable}
